\exercisegroup

¶ 1) 设 E 是一个可度量的局部凸空间，\(E_{b}^{\prime}\) 是它的强对偶。若 \(E_{b}^{\prime}\) 是可度量的，证明 E 的拓扑可以由单独一个范数定义（利用 III, 第 37 页, exerc. 2 与第 38 页, exerc. 5，以及 E 是有界型的这一事实）。

\begin{enumerate}
\def\labelenumi{\arabic{enumi})}
\setcounter{enumi}{1}
\tightlist
\item
  次桶型空间是半桶型的。我们称一个局部凸空间为 (DF) 空间，如果它是半桶型的，并且其典范有界结构 (III, 第 3 页, 定义 5) 具有可数基。每个赋范空间以及赋范空间序列的每个严格归纳极限 (II, 第 33 页) 都是 (DF) 空间。Fréchet 空间的每个强对偶都是 (DF) 空间。
\end{enumerate}

\begin{enumerate}
\def\labelenumi{\alph{enumi})}
\item
  (DF) 空间的强对偶是 Fréchet 空间。
\item
  设 \(\mathrm{E}\) 是一个 (DF) 空间，\((\mathbf{A}_n)\) 是 \(\mathrm{E}\) 中有界、凸、均衡且闭的子集的一个递增序列，使得 \(\mathrm{E}\) 的每个有界子集都被诸 \(\mathbf{A}_n\) 之一吸收。
\end{enumerate}

  设 U 是诸 \(A_{n}\) 的并；证明：U 在 E 中的闭包 \(\overline{U}\) 恰好是所有这样的 \(x \in E\) 组成的集：\(\lambda x \in U\) 对 \(0 \leqslant \lambda < 1\) 。（若 \(x \notin \lambda U\) 对某个 \(\lambda > 1\) ，则对每个 n，存在线性型 \(x_{n}' \in E'\) 使得 \(x_{n}' \in A_{n}^{\circ}\) 且 \(\langle x, x_{n}' \rangle = \lambda\) ，而序列 \((x_{n}')\) 是等度连续的，从而有一个弱极限点。）

\begin{enumerate}
\def\labelenumi{\alph{enumi})}
\setcounter{enumi}{2}
\tightlist
\item
  证明：若一个 (DF) 空间是桶型的，则它也是有界型的（参见 III, 第 44 页, exerc. 13, b)）。给出非超有界型但是有界型且桶型的 (DF) 空间的例子 (III, 第 46 页, exerc. 22)，以及非桶型但是有界型的 (DF) 空间的例子。
\end{enumerate}

\begin{enumerate}
\def\labelenumi{\arabic{enumi})}
\setcounter{enumi}{2}
\tightlist
\item
  设 \(E\) 是一个可度量的局部凸空间，\(E_b'\) 是它的强对偶。
\end{enumerate}

\begin{enumerate}
\def\labelenumi{\alph{enumi})}
\item
  证明：每个凸均衡子集 \(V'\) （属于 \(E_{b}'\) ）只要吸收 \(E_{b}'\) 的强有界子集，就包含一个（对于强拓扑的）桶，该桶吸收 \(E_{b}'\) 的强有界子集。（设 \((\mathbf{K}_{n}')\) 是 \(E_{b}'\) 的典范有界结构的一个可数基，并设 \(\lambda_{n}\) 使得 \(\lambda_{n}K_{n}' \subset \frac{1}{2}V'\) ；把 exerc. 2, b) 应用于序列 \(A_{n}'\) ，其中 \(A_{n}'\) 是诸 \(\lambda_{j}K_{j}'\) （\(j \leqslant n\) ）的并的凸包。）
\item
  由 \(a\) ) 推出下列性质等价：
\end{enumerate}

α) E 是区分的 (IV, 第 52 页, exerc. 4)。

β) \(E_{b}^{\prime}\) 是次桶型的 (III, 第 44 页, exerc. 7)。

\(\gamma)\) \(\mathrm{E}_b^{\prime}\) 是有界型的。

δ) \(E_{b}^{\prime}\) 是桶型的。

ε) \(E_{b}^{\prime}\) 是超有界型的 (III, 第 45 页, exerc. 19)。

\begin{enumerate}
\def\labelenumi{\alph{enumi})}
\setcounter{enumi}{2}
\tightlist
\item
  证明：若 \(E_{b}^{\prime}\) 是自反的，则 \(\hat{E} = E''\) （它显然是自反的）（参见 IV, 第 52 页, exerc. 4 及第 53 页, exerc. 11）。
\end{enumerate}

\begin{enumerate}
\def\labelenumi{\arabic{enumi})}
\setcounter{enumi}{3}
\tightlist
\item
  设 E 是一个 Hausdorff 局部凸空间，\(E'\) 是它的对偶。若 M 是 E 的一个可度量且区分的 (IV, 第 52 页, exerc. 4) 闭向量子空间，则强拓扑 \(\beta(\mathrm{E}'/\mathrm{M}^{\circ}, \mathrm{M})\) 是按 \(M^{\circ}\) 从强拓扑 \(\beta(\mathrm{E}', \mathrm{E})\) 得到的商拓扑（利用 exerc. 3, b) 与 IV, 第 51 页, exerc. 22, b)）。
\end{enumerate}

¶ 5) 对每个整数 n \textgreater{} 0，设 \(a^{(n)}\) 是二重序列 \((a_{pq}^{(n)})\) （ \(p \in N\) ，\(q \in N\) ），满足：\(a_{pq}^{(n)} = q\) 若 \(p \leqslant n\) ，且 \(a_{pq}^{(n)} = 1\) 若 p \textgreater{} n。~设 E 是由所有满足如下条件的实数二重序列 \(x = (x_{pq})_{(p,q) \in \mathbf{N} \times \mathbf{N}}\) 组成的向量空间：对每个整数 n \textgreater{} 0，数 \(r_n(x) = \sum_{p,q} a_{pq}^{(n)} |x_{pq}|\)

是有限的。若给 E 赋予由诸半范数 \(r_n\) 定义的拓扑，则 E 是一个满足第一可数公理的 Fréchet 空间 (IV, 第 47 页, exerc. 1, c))；E 的对偶 E' 可以与由所有满足如下条件的实数二重序列 \(x' = (x_{pq}')\) 组成的空间等同：至少对一个指标 \(n\) ，存在 \(k_n > 0\) 使得 \(|x_{pq}'| \leqslant k_n a_{pq}^{(n)}\) 对每一对 \((p, q)\) ；并且 \(\langle x, x' \rangle = \sum_{p, q} x_{pq} x_{pq}'\) 。

(IV, 第 47 页, exerc. 1, c))。

对每个整数 \(p_0 > 0\) 以及整数的每个序列 \((m_p)\) （诸整数 \(> 0\) ），设 \(\mathrm{J}(p_0; (m_p))\) 是由所有整数对 \(p > 0\) ，\(q > 0\) （满足 \(p \geqslant p_0\) 且 \(q \geqslant m_p\) ）组成的集；设 \(\mathfrak{V}\) 是 \(\mathbf{N} \times \mathbf{N}\) 上由诸集 \(\mathrm{J}(p_0; (m_p))\) 组成的滤子基，并设 \(\mathfrak{F}\) 是一个比以 \(\mathfrak{V}\) 为基的滤子更细的超滤子。

\begin{enumerate}
\def\labelenumi{\alph{enumi})}
\item
  证明：对所有 \(x' = (x_{pq}') \in \mathrm{E}'\) ，二重序列 \((x_{pq}')\) 有一个极限 \(u(x')\) （关于超滤子 \(\mathfrak{F}\) ）；若 \(\mathrm{V}_n\) 是 \(\mathrm{E}\) 中由 \(r_n(x) \leqslant 1\) 定义的 0 的邻域，则 \(|u(x')| \leqslant 1\) 对所有 \(x' \in \mathrm{V}_n^\circ\) 。
\item
  设 \(U'\) 是 \(E'\) 中对于强拓扑的 0 的一个邻域，它是凸的、均衡的且弱闭的，并对每个 n，设 \(\alpha_{n} > 0\) 使得 \(\alpha_{n} V_{n}^{\circ} \subset U'\) 。对每个整数 p \textgreater{} 0，设 \(m_{p}\) 是使得 \(2^{p+1} \leqslant \alpha_{p} m_{p}\) 的整数，并设 \(x' = (x'_{pq})\) 是满足如下条件的二重序列：\(x'_{pq} = 0\) 对 q \textless{} \(m_{p}\) ，\(x'_{pq} = 2\) 对 \(q \geqslant m_{p}\) 。证明：\(x' \in U'\) 但 \(u(x') = 2\) ；由此推出 u 在 \(E'\) 中不是强连续的，尽管它在 \(E'\) 的每个有界子集上是有界的。推断 (IV, 第 58 页, exerc. 3)：E 不是区分的，从而强对偶 \(E'_{b}\) 是一个非次桶型的 (DF) 空间。
\item
  利用 b) 构造一个 Fréchet 空间 F 的闭子空间 M 的例子，使得强拓扑 \(\beta(\mathrm{F}^{\prime}/\mathrm{M}^{\circ}, \mathrm{M})\) 不同于按 \(M^{\circ}\) 从强拓扑 \(\beta(\mathrm{F}^{\prime}, \mathrm{F})\) 得到的商拓扑（把 E 嵌入到一个 Banach 空间的可数乘积中）。
\end{enumerate}

¶ 6) a) 设 E 是一个 (DF) 空间 (IV, 第 57 页, exerc. 2)，并设 U 是一个凸集，使得对 E 的每个有界、凸且均衡的子集 A，U ∩ A 都是 0 的一个邻域（对于 E 的拓扑在 A 上诱导的拓扑）。证明：U 是 A 中 0 的一个邻域。（设 (A \(_{n}\) ) 是 E 的典范有界结构的一个可数基 (III, 第 3 页, 定义 5)。证明：可以用归纳法定义一个序列 ( \(\lambda_{n}\) ) （由数 \textgreater{} 0 组成）以及一个序列 (V \(_{n}\) ) （由闭的、凸且均衡的

0 的邻域组成），使得 \(\lambda_{n}\mathrm{A}_{n} \subset \frac{1}{3}\mathrm{U}\) ，\(\lambda_{n}\mathrm{A}_{n} \subset \bigcap_{j=1}^{\infty} \mathrm{V}_{j}\) ，\(\mathrm{V}_{n} \cap \mathrm{A}_{n} \subset \mathrm{U}\) 对每个 \(n\) 。

首先证明：若 \(\lambda_{j}\) 与 \(\mathrm{V}_j\) 已对 \(j\leqslant n\) 构造出来，则可找到 \(\lambda_{n + 1}\) 使得 \(\lambda_{n + 1}\mathrm{A}_{n + 1}\subset \frac{1}{3}\mathrm{U}\) 且 \(\lambda_{n + 1}\mathrm{A}_{n + 1}\subset \mathrm{V}_j\) 对 \(j\leqslant n\) 。其次证明：可找到 \(\mathrm{V}_{n + 1}\) 使得 \(\lambda_j\mathrm{A}_j\subset \mathrm{V}_{n + 1}\) 对 \(j\leqslant n + 1\) 且 \(\mathrm{V}_{n + 1}\cap \mathrm{A}_{n + 1}\subset \mathrm{U}\) ；为此，设 A 表示诸 \(\lambda_{j}\mathrm{A}_{j}\) （\(j\leqslant n + 1\) ）的凸包，证明可取 \(\mathrm{V}_{n + 1} = \overline{\mathrm{A} + \mathrm{V}}\) ，其中 V 是一个适当的凸均衡 0 邻域；我们指出，为此只需证明：若 \(\mathbf{B} = \mathbf{A}_{n + 1}\cap \mathbb{C}\) U，则 0 不在集 \(\mathbf{B} + 2\mathbf{A}\) 的闭包中。

\begin{enumerate}
\def\labelenumi{\alph{enumi})}
\setcounter{enumi}{1}
\tightlist
\item
  由 a) 推出：若 u 是从 E 到某个局部凸空间 F 中的线性映射，且 u 在 E 的每个有界子集上的限制都是连续的，则 u 是连续的（参见 IV, 第 50 页, exerc. 15）。
\end{enumerate}

¶ 7) a) 设 E 是一个 (DF) 空间，U 是 E 中一个凸、均衡且闭的集，它吸收 E 的有界子集，并设 \((x_n)\) 是 \(\complement\) U 的点的一个序列。证明：存在 E 中 0 的一个邻域 V，它不包含任何一个 \(x_n\) 。（设 \((A_n)\) 是 E 的典范有界结构的一个可数基。证明：可以用归纳法定义一个序列 \((\lambda_n)\) （由数 \textgreater{} 0 组成）以及一个序列 \((V_n)\) （由凸、均衡且闭的

0 的邻域组成），使得 \(\lambda_{n}\mathbf{A}_{n}\subset \bigcap_{j = 1}\mathbf{V}_{j},\lambda_{n}\mathbf{A}_{n}\subset \mathrm{U}\) 且 \(x_{n}\in \mathbb{C}\mathbf{V}_{n}\) 对所有 \(n\) 。为此，若诸 \(\lambda_{j}\) 与 \(\mathrm{V}_j\) 已对

\(j \leqslant n\) 构造出来，则取 \(\lambda_{n+1}\) 使得 \(\lambda_{n+1}\mathrm{A}_{n+1} \subset \mathrm{U}\) 且 \(\lambda_{n+1}\mathrm{A}_{n+1} \subset \mathrm{V}_j\) 对所有 \(j \leqslant n\) ，然后取 \(\mathrm{V}_{n+1}\) 包含诸 \(\lambda_j\mathrm{A}_j\) （\(j \leqslant n + 1\) ）的并的凸包的闭包。

\begin{enumerate}
\def\labelenumi{\alph{enumi})}
\setcounter{enumi}{1}
\item
  由 a) 推出：若 M 是 E 的一个包含处处稠密可数集的子集，则 E 的双对偶 \(E''\) 的强拓扑在 M 上诱导的拓扑与 E 的拓扑在 M 上诱导的拓扑相同。特别地，E 中的收敛序列对于 E 的拓扑与对于 \(E''\) 的强拓扑诱导的拓扑是同样的；对 E 的每个可度量子集 M，E 的拓扑在 M 上诱导的拓扑与 \(E''\) 的强拓扑诱导的拓扑相同。
\item
  由 \(a\) ) 推出：若在 \(E\) 中存在一个可数的处处稠密集，则 \(E\) 是次桶型的。
\item
  由 b) 以及 exerc. 6 推出：若 E 的每个有界子集对于 E 的拓扑诱导的拓扑都是可度量的，则 E 是次桶型的。
\end{enumerate}

\begin{enumerate}
\def\labelenumi{\arabic{enumi})}
\setcounter{enumi}{7}
\tightlist
\item
  设 E 是一个 Fréchet 空间，\(E_{b}^{\prime}\) 是它的强对偶。假设存在一个处处稠密序列 \((x_{n}^{\prime})\) 于 \(E_{b}^{\prime}\) 中。证明：E 满足第一可数公理。（设 \((\mathbf{K}_{n}^{\prime})\) 是 \(E_{b}^{\prime}\) 的典范有界结构的一个由闭凸均衡集组成的可数基。对每个系统 \(\alpha\) ——它由一点 \(x_{n}^{\prime}\) 、任意有限个有理数 \(\lambda_{k} > 0 (1 \leqslant k \leqslant m)\) 以及 m 个指标 \(n_{k}\) 组成，且满足 \(x_{n}^{\prime} \notin 2 \sum_{k=1}^{m} \lambda_{k} K_{n_{k}}^{\prime} = 2 H_{\alpha}^{\prime}\) ——取 \(x_{\alpha} \in E\) 使得
\end{enumerate}

方程为 \(\langle x_{\alpha}, y' \rangle = 1\) 的超平面把两个弱紧集 \(\mathrm{H}_{\alpha}'\) 与 \(x_{n}' + \mathrm{H}_{\alpha}'\) 严格分离。证明：对每个 \(x' \neq 0\) （属于 \(\mathrm{E}'\) ），存在一个系统 \(\alpha\) 使得 \(\langle x_{\alpha}, x' \rangle \neq 0\) 。为此，考虑 0 的一个邻域 \(\mathrm{V}'\) （\(\mathrm{E}_b'\) 中的），使得 \(\mathrm{V}' \cap (x' + \mathrm{V}') = \emptyset\) ，然后对每个整数 \(m\) ，取有理数 \(\lambda_m > 0\) 使得 \(\lambda_m \mathrm{K}_m' \subset \mathrm{V}'\) ；利用如下事实：并 \(\mathrm{U}' \subset \mathrm{V}'\) （由诸 \(\lambda_m \mathrm{K}_m'\) 取并）是 0 的一个邻域（exerc. 7，以及 IV, 第 58 页, exerc. 3, b)），并且存在 \(n\) 使得 \(x_n' \in x' + \mathrm{U}'\) 。

¶ 9) a) 设 E 是一个 Hausdorff 半桶型空间，M 是 E 的一个闭向量子空间，E' 是 E 的对偶。证明：E/M 是半桶型的，并且强拓扑 β(M°, E/M) 与强拓扑 β(E', E) 在 M° 上诱导的拓扑相同。（注意，只需证明 M° 中一个对于 β(E', E) 收敛于 0 的序列 (x') 对于 β(M°, E/M) 是有界的。）由此推出：若 E 是 (DF) 空间，则 E/M 也是（参见 IV, 第 63 页, exerc. 8）。

\begin{enumerate}
\def\labelenumi{\alph{enumi})}
\setcounter{enumi}{1}
\item
  设 E 是一个 Hausdorff 局部凸空间，M 是 E 的一个（未必闭的）向量子空间。证明：若 M 是半桶型空间，则强拓扑 \(\beta(\mathrm{E}^{\prime}/\mathrm{M}^{\circ}, \mathrm{M})\) 与按 \(M^{\circ}\) 从强拓扑 \(\beta(\mathrm{E}^{\prime}, \mathrm{E})\) 得到的商拓扑相同。（按 a) 的方式论证。）
\item
  证明：一个 Hausdorff 且拟完备的半桶型空间 E 是完备的（对 E 与 \(\hat{E}\) 应用 b)）。特别地，一个半桶型、半自反的空间是完备的（参见 IV, 第 52 页, exerc. 6）。
\item
  证明：半桶型（相应地，(DF)）Hausdorff 空间的完备化是半桶型的（相应地，是 (DF) 空间）。
\item
  设 \((\mathrm{E}_n)\) 是半桶型（相应地，(DF)）空间的一个序列，E 是一个向量空间，且对每个 \(n\) ，设 \(f_n\) 是从 \(\mathrm{E}_n\) 到 E 中的一个线性映射。假设 \(\dot{\mathrm{E}}\) 是诸 \(f_n(\mathrm{E}_n)\) 的并；证明：E 对于使诸 \(f_n\) 都连续的最细局部凸拓扑是半桶型的（相应地，是 (DF) 空间）（先考察 E 是诸 \(\mathrm{E}_n\) 的拓扑直和的情形）。若诸 \(\mathrm{E}_n\) 都是半自反的（相应地，自反的）且 E 是 Hausdorff 的，则 E 是半自反的（相应地，自反的）。
\end{enumerate}

\begin{enumerate}
\def\labelenumi{\arabic{enumi})}
\setcounter{enumi}{9}
\tightlist
\item
  设 E 是一个满足第一可数公理的 Fréchet 空间。证明：若在 E 的对偶 \(E'\) 中，每个对于弱拓扑 \(\sigma(E', E)\) 收敛的序列也对于强拓扑 \(\beta(E', E)\) 收敛，则 E 是一个 Montel 空间。（证明：\(E'\) 的每个有界子集对于强拓扑都是相对紧的；利用 GT, II, § 4, exerc. 6；然后利用 IV, 第 53 页, exerc. 11, b)。）
\end{enumerate}

¶ 11) 设 \((c_{mn})\) 是由 \textgreater{} 0 的数组成的一个二重序列，满足 \(c_{m,n} \leqslant c_{m+1,n}\) ，并设 E 是由所有实数序列 \(x = (x_n)\) 组成的空间，这些序列使得 \(p_m(x) = \sum_n c_{mn} |x_n| < +\infty\) ，对

每个整数 \(m\) 成立。我们给 \(E\) 赋予由诸半范数 \(p_m\) 定义的拓扑，则对于该拓扑 \(E\) 是一个满足第一可数公理的 Fréchet 空间；对偶 \(E'\) （\(E\) 的对偶）可以与由所有满足如下条件的序列 \(x' = (x_n')\) 组成的空间等同：\(\sup_{n} c_{mn}^{-1} |x_n'| < +\infty\) 至少对

一个 \(m\) 成立，且典范双线性型 \(\langle x, x' \rangle\) 等同于 \(\sum_{n} x_n x_n'\) (IV, 第 47 页,

exerc. 1, c))。假设不存在任何子序列 \((n_k)\) ，使得有一个序列 \((a_m)\) （由 \(\geqslant 0\) 的数组成）以及一个指标 \(m_0\) ，满足 \(c_{m,n_k} \leqslant a_m c_{m_0,n_k}\) 对所有 \(m \geqslant m_0\) 以及所有 \(k\) 。在这些条件下，\(\mathbf{E}'\) 中每个弱收敛序列都强收敛，从而 (IV, 第 60 页, exerc. 10) E 是一个 Montel 空间。（用反证法；若有必要，作一个形如 \((x_n) \mapsto (a_n x_n)\) 的变换，化归为 \(c_{m_0 n} = 1\) 对所有 \(n\) 及某个 \(m_0\) 成立的情形，且存在一个序列 \((x'^{(p)})_{p \geqslant 0}\) 在 \(\mathbf{E}'\) 中弱收敛于 0，并满足 \(|x_n'^{(p)}| \leqslant 1\) 对每一对 \((p, n)\) ，同时存在 E 中一个有界集 B，由 \(p_m(x) \leqslant b_m\) （对所有 \(m \geqslant 0\) ）定义，且使得 \(\sup_{x \in \mathbf{B}} |\langle x, x'^{(p)} \rangle| \geqslant 2\delta > 0\) 对

每个整数 \(p\) 成立。在这些假设下，证明存在整数的一个严格递增序列 \((r_q)\) ，以及点的一个序列 \((x^{(q)})\) （\(\mathbf{B}\) 的点），使得

\[
\sum_ {k = r _ {q} + 1} ^ {r _ {q+1}} \left| x _ {k} ^ {(q)} \right| > \delta
\]

对每个指标 \(q\) 成立。然后（用反证法）证明：对每个 \(q\) ，至少存在一个指标 \(s_q\) 使得 \(r_q < s_q \leqslant r_{q+1}\) ，并且对每个整数 \(m\) 有 \(c_{m,s_q} \leqslant b_m 2^{m+2}/\delta\) ，这与假设矛盾。）

\begin{enumerate}
\def\labelenumi{\arabic{enumi})}
\setcounter{enumi}{11}
\tightlist
\item
  \begin{enumerate}
  \def\labelenumii{\alph{enumii})}
  \tightlist
  \item
    设 F 是一个 Hausdorff (DF) 空间，\(F_{b}^{\prime}\) 是它的强对偶。证明：若 \(F_{b}^{\prime}\) 是自反的，则 F 的完备化 \(\hat{F}\) 是自反的且等于 F 的双对偶 \(F^{\prime\prime}\) （参见 IV, 第 52 页, exerc. 4 及第 53 页, exerc. 11）。
  \end{enumerate}
\end{enumerate}

\begin{enumerate}
\def\labelenumi{\alph{enumi})}
\setcounter{enumi}{1}
\tightlist
\item
  设 E 是一个 Fréchet 空间。证明：若 E 的双对偶 \(E''\) 是自反的，则 E 是自反的。
\end{enumerate}

\begin{enumerate}
\def\labelenumi{\arabic{enumi})}
\setcounter{enumi}{12}
\tightlist
\item
  \begin{enumerate}
  \def\labelenumii{\alph{enumii})}
  \tightlist
  \item
    设 E、F 是两个 Fréchet 空间，G 是一个 Hausdorff 局部凸空间，\(E'\) ，\(F'\) ，\(G'\) 分别是 E、F、G 的对偶。设 u 是从 \(E' \times F'\) 到 \(G'\) 中的一个双线性映射，当 \(E'\) ，\(F'\) ，\(G'\) 分别带上弱拓扑 \(\sigma(E', E)\) ，\(\sigma(F', F)\) 与 \(\sigma(G', G)\) 时，u 是分别连续的 (III, 第 28 页) 。证明：在这些条件下，u 是从 \(E' \times F'\) 到 \(G'\) 中的一个连续映射（当 \(E'\) ，\(F'\) 与 \(G'\) 带上强拓扑时）。（对 \(z \in G\) ，令 \(\langle z, u(x', y') \rangle = \langle v_z(x'), y' \rangle\) ，其中 \(v_z(x') \in F\) 。首先证明：若 \(E'\) 带上强拓扑而 \(F\) 带上初始拓扑，则诸 \(v_z\) （当 \(z\) 取遍一个有界集 \(C\) （\(G\) 中的）时）所成的集是等度连续的；为此利用 IV, 第 51 页, exerc. 19, d)。其次证明：存在 0 的一个邻域 \(V'\) （对于 \(E'\) 的强拓扑），使得诸集 \(v_z(V')\) （当 \(z\) 取遍 \(C\) 时）的并在 \(F\) 中是有界的；为此利用 III, 第 47 页, exerc. 5。）
  \end{enumerate}
\end{enumerate}

\begin{enumerate}
\def\labelenumi{\alph{enumi})}
\setcounter{enumi}{1}
\tightlist
\item
  给出一个例子，以说明若假设 E 是 Fréchet 空间而 F 是 Fréchet 空间的一个严格归纳极限 (III, 第 47 页, exerc. 3)，则 a) 的结论不成立。
\end{enumerate}

\begin{enumerate}
\def\labelenumi{\arabic{enumi})}
\setcounter{enumi}{13}
\tightlist
\item
  \begin{enumerate}
  \def\labelenumii{\alph{enumii})}
  \tightlist
  \item
    设 E 是一个 Fréchet 空间，\(E'\) 是它的对偶。证明：\(E'\) 带上紧收敛拓扑或任何更细的 S-拓扑都是完备的（参见 III, 第 22 页, 注记 1）。若 E 不是自反的，证明：\(E'\) 对于任何比紧收敛拓扑更细且比 \(\tau(\mathrm{E}', \mathrm{E})\) 更粗的 S-拓扑都不是次桶型的。
  \end{enumerate}
\end{enumerate}

\begin{enumerate}
\def\labelenumi{\alph{enumi})}
\setcounter{enumi}{1}
\tightlist
\item
  设 \((\mathrm{E}_{\alpha})_{\alpha\in\mathbf{A}}\) 是一族 Fréchet 空间，E 是一个向量空间，且对每个 \(\alpha\in A\) ，设 \(h_{\alpha}\) 是从 \(E_{\alpha}\) 到 E 中的一个线性映射。假设 E 带上使诸 \(h_{\alpha}\) 都连续的最细局部凸拓扑 (II, 第 27 页) 后是 Hausdorff 的。证明：E 的对偶 \(E'\) 带上紧收敛拓扑或任何更细的 S-拓扑都是完备的（参见 III, 第 20 页, 定理 1）。
\end{enumerate}

\begin{enumerate}
\def\labelenumi{\arabic{enumi})}
\setcounter{enumi}{14}
\tightlist
\item
  设 \(\mathrm{E}\) 是一个无限维 Banach 空间，\((a_{n})_{n\geqslant 1}\) 是 \(\mathrm{E}\) 的点的一个可数、完全且自由的族，并设 \(\mathrm{F}_n\) 是一个 \(n\) 维子空间（\(\mathrm{E}\) 中的），由诸 \(a_{m}\) （\(m\leqslant n\) ）生成。设 \(\mathrm{S}_n\) 是方程为 \(\| x\| = n\) 的球面（\(\mathrm{E}\) 中的）；在 \(\mathrm{S}_n\cap \mathrm{F}_n\) 中设 \(\mathrm{A}_n\) 是一个有限集，使得
\end{enumerate}

\(S_{n} \cap F_{n}\) 的每一点都有距它 \(\leqslant 1/n\) 的 \(A_{n}\) 中的点。证明：\(A = \bigcap_{n=1}^{\infty} A_{n}\) 与每个闭有界集的交

都是闭的，但 0 是 A 对于弱化拓扑的一个极限点。

\begin{enumerate}
\def\labelenumi{\arabic{enumi})}
\setcounter{enumi}{15}
\tightlist
\item
  设 E 是局部凸可度量空间的一个序列 \((\mathrm{E}_{p})\) 的归纳极限空间，典范映射 \(E_{p} \rightarrow E\) 都是单射，且 E 是诸 \(E_{p}\) 的像的并。证明：E 的强对偶 \(E_{b}^{\prime}\) 是可穷竭的 (III, 第 49 页, exerc. 1)。（设 \((\mathrm{U}_{j}^{p})\) 是 \(E_{p}\) 中 0 的闭、凸且均衡邻域的一个递减基本系；考虑极集 \((\mathrm{U}_{j}^{p})^{\circ}\) （\(E^{\prime}\) 中的）的有限交；利用如下事实：对每个递增序列
\end{enumerate}

\((m_{j})_{j\geqslant 1}\) ，交 \(\bigcap_{j = 1}^{\infty}(\mathrm{U}_{m_j}^j)^{\circ}\) 是 E 中 0 的某个邻域的极集，并且 \(\mathrm{E}_b^\prime\) 是完备的。）

¶ 17) a) 设 E 是一个 Hausdorff 局部凸空间，使得由 E 的相对紧集组成的有界结构具有一个可数基 (A \(_{n}\) ) (III, 第 1 页)。证明：(A \(_{n}\) ) 也是典范有界结构的一个基 (III, 第 3 页, 定义 5)。（设 C \(_{n}\) 是 n 个都等于 A \(_{n}\) 的集之和所成的相对紧集；则 (C \(_{n}\) ) 也是相对紧集有界结构的一个基。用反证法，考虑一个不包含于任何 C \(_{n}\) 中的有界集 B，推断存在 B 的点的一个序列 (x \(_{n}\) ) 使得 x \(_{n}\) /n∉ A \(_{n}\) ；由此导出矛盾。）于是，空间 E 是半自反的，并且 E 中每个紧集的闭凸均衡包都是紧的。

\begin{enumerate}
\def\labelenumi{\alph{enumi})}
\setcounter{enumi}{1}
\tightlist
\item
  假设 E 是次桶型的，并且对一个与 E 和 \(E'\) 之间的对偶相容的拓扑 T，E 中对于 T 的相对紧子集组成的有界结构具有可数基。证明：此时 E 是一个自反的 (DF) 空间 (IV, 第 57 页, exerc. 2)。（利用 a)，注意 E 中的闭有界集对于 T 都是紧的，从而对于 E 的初始拓扑是完备的；这就蕴含 E 是桶型的。）若 T 是初始拓扑，则 E 是一个 Montel 空间。
\end{enumerate}

\begin{enumerate}
\def\labelenumi{\arabic{enumi})}
\setcounter{enumi}{17}
\tightlist
\item
  \begin{enumerate}
  \def\labelenumii{\alph{enumii})}
  \tightlist
  \item
    设 E 是一个 Hausdorff 局部凸空间，\((\mathbf{A}_{n})\) 是对于弱化拓扑 \(\sigma(\mathrm{E}, \mathrm{E}')\) 凸、均衡、紧的集的一个递增序列，满足诸集 \(A_{n}\) 的并是 E，且对每个整数 n 及每个 \(\lambda > 0\) ，存在 m 使得 \(\lambda A_{n} \subset A_{m}\) 。证明：\((\mathbf{A}_{n})\) 是由对于 \(\sigma(\mathrm{E}, \mathrm{E}')\) 凸且相对紧的集组成的有界结构的一个基。（注意：在 \(E'\) 上，在诸 \(A_{n}\) 上一致收敛的拓扑就是 \(\tau(\mathrm{E}', \mathrm{E})\) 。）b) 若 E 是桶型的，且存在一个具有 a) 中所述性质的序列 \((\mathbf{A}_{n})\) ，则 E 是一个自反的 (DF) 空间。（注意：\(E'\) 带上 \(\tau(\mathrm{E}', \mathrm{E})\) 是可度量的，并且 E 的拓扑是 \(\beta(\mathrm{E}, \mathrm{E}')\) 。）
  \end{enumerate}
\end{enumerate}

